% =====================================================================
%  Chapter 28 opener -- "Chapter at a Glance", three reading paths + the
%  ladder staircase, THE running example (the 9-node rooted tree R whose
%  root has two isomorphic size-3 children plus one size-2 child, its
%  unrooted twin R', the
%  look-alike pair (T1, T2) that shares every cheap invariant yet is not
%  isomorphic, and the two-centre path P6), the
%  who-owns-what comparison against Ch.12 (centre/diameter), Ch.13 (tree DP),
%  Ch.17 (small-to-large), Ch.15 (rerooting), Ch.27 (Mo boundary) and
%  Ch.26 (virtual trees), the recognition guide, a constraint->technique
%  table, and notation.
%  Exports fig:c22-tree, fig:c22-twin, fig:c22-trap, fig:c22-taxonomy.
%  No \term{centre} / \term{diameter} here -- they are indexed by Ch.12;
%  \term the NEW objects "isomorphic", "canonical label", "interned tuple",
%  "subtree hash", "commutative combine", "symmetrical" only.
% =====================================================================

\begin{chaptertheorybox}{Chapter at a Glance}
Chapter~13 (\chref{c10}) taught you to fold \emph{numbers} bottom-up: a value at
each node computed from its children's values.  Chapter~12 (\chref{c09}) taught
you to find a tree's \emph{centre} -- the one or two vertices minimizing the
maximum distance -- via the diameter.  Both are tools for measuring a tree.
This chapter is about \emph{comparing} trees.

Two trees are \term{isomorphic} when one is a relabelling of the other: a
bijection of vertices preserving adjacency.  Labels vanish, child order vanishes
-- only the \textbf{shape} remains, and shape equality is the question behind a
surprising stack of problems: ``is this subtree the same shape as that one?''
(Ch.27's boundary parting question, \chref{c21}), ``can the children be ordered
to make a mirror image?'' (CF 1800G), ``cut one node -- do the pieces all look
alike?'' (CF 1252F).  There is no known fast general algorithm for graph
isomorphism, but \emph{trees} are exactly where the problem becomes easy: the
\term{canonical label} of an unordered rooted tree is computed by one bottom-up
pass that sorts child labels and \term{interns each sorted tuple} to a fresh id
(the AHU algorithm), and an unrooted tree is canonized by rooting it at its
centre.  The randomized cousin -- a \term{subtree hash} built with a
\term{commutative combine} -- answers ``same shape?'' in \(O(1)\) per compare
after \(O(n)\) preprocessing, at the price of an explicit collision ledger.

\textbf{You will learn to}
\begin{itemize}[nosep]
  \item run the AHU fold by hand: leaves first, sort each node's child labels,
        map each sorted tuple to a fresh id; two rooted trees are isomorphic iff
        their root ids match (\(\secref{c22-ahu}\)); \hfill\lv{3}
  \item hash subtrees with a commutative combine (depth-seeded XOR/sum of
        per-child values), size the collision risk, and re-verify a claimed YES
        with an explicit bijection DFS (\(\secref{c22-hash}\), solved CF 601D
        \emph{Acyclic Organic Compounds}); \hfill\lv{4}--\lv{5}
  \item decide \term{unrooted} isomorphism: root each tree at its centre (the
        Ch.12 black box), compare canonical forms, and handle the two-centre
        edge case by pairing halves (\(\secref{c22-centre}\), solved CF 1800G
        \emph{Symmetree}); \hfill\lv{4}
  \item count distinct subtree \emph{shapes} by interned ids -- isomorphism
        classes and their multiplicities in \(O(n\log n)\)
        (\(\secref{c22-distinct}\)); \hfill\lv{4}
  \item read a statement and pick canonical-vs-hash vs plain tree DP, and know
        exactly when shape equality is a red herring
        (\(\secref{c22-boundary}\)). \hfill\lv{5}
\end{itemize}

\textbf{Prerequisites.} Chapter~12 (\chref{c09}) -- diameter and centre: the
canonical root, used as a black box.  Chapter~13 (\chref{c10}) -- the bottom-up
fold, reused here with \emph{shapes} as the folded object.  Chapter~17
(\chref{c28}) -- small-to-large set merging: the engine behind the 601D solve.
Chapter~15 (\chref{c20}) -- rerooting: parent-side hashes in the Summit.
Chapter~27 (\chref{c21}) -- the \textbf{previous chapter}; its method boundary
hands the ``which subtree \emph{equals} which'' question to this one.
\end{chaptertheorybox}

% ---------------- Reading paths + ladder position ----------------
\begin{chaptertheorybox}{How to Read This Chapter (three paths)}
Level badges \(\lv3\)--\lv{5}: \(\lv3\) is running the AHU labelling by hand on
small trees and spotting the wrong-root trap; \(\lv4\) is the two coded ideas --
centre-rooted comparison (1800G) and interned-id shape counting; \(\lv{5}\) is
the hash-set small-to-large merge (601D) and the all-components compare of the
Summit (CF 1252F).

\textbf{Deterministic path (\(\lv3\)--\(\lv4\), ``never lose to a collision'').}
\secref{c22-ahu} in full -- sort, intern, compare ids -- then \secref{c22-centre}
for the unrooted lift (centre as canonical root, solved CF 1800G), capped by
\secref{c22-distinct}: counting shape classes with the same interned ids.

\textbf{Randomized path (\(\lv4\)--\lv{5}, ``\(O(1)\) compares'').}
\secref{c22-hash} -- the commutative combine, 64-bit moduli, the collision
ledger -- then the solved CF 601D, where per-node sets of downward-string hashes
merge small-to-large; the Summit \(\probref{c22-summit}\) compares hashed
components of every candidate cut.

\textbf{Boundary path (\(\lv4\)--\lv{5}, ``is this even a shape job?'').}
\secref{c22-boundary}: the tell (labels and child order are irrelevant, the
question is \emph{equality of unlabelled structures}) and the look-alikes that
are not -- labelled path strings (suffix tools), per-node value aggregates
(tree DP), online edits (HLD/link-cut), ancestor set queries (Ch.27's Mo).
\end{chaptertheorybox}

\begin{figure}[H]
\centering
\begin{tikzpicture}[diagram, x=8.5mm, y=6.5mm]
  \foreach \st/\lab in {1/lv1,2/lv2,3/lv3,4/lv4,5/lv5,6/lv6} {
    \fill[\lab!22, draw=\lab, line width=0.8pt] (\st-1,0) rectangle (\st,\st*0.52);
    \node[font=\scriptsize\bfseries, text=\lab] at (\st-0.5,\st*0.52-0.42) {L\st};}
  \node[annot, anchor=south west, font=\scriptsize]
    at (0.05,0.62) {prereq (Ch.12/13): centre black box, the bottom-up fold};
  \node[annot, anchor=south west, font=\scriptsize]
    at (2.05,1.66) {28.1 sort child labels, intern tuples -- the AHU fold by hand};
  \node[annot, anchor=south west, font=\scriptsize]
    at (3.05,2.18) {28.3 solved CF 1800G (mirror pairing) | 28.4 shape counting};
  \node[annot, anchor=south west, font=\scriptsize, text=lv5]
    at (4.05,2.70) {28.2 solved CF 601D (hash sets + small-to-large)};
  \node[annot, anchor=south west, font=\scriptsize, text=lv6]
    at (5.05,3.22) {28.5 the pick; Summit CF 1252F (all components identical?)};
\end{tikzpicture}
\caption{Chapter~28 on the six-step ladder: hand-running the AHU fold and
spotting the wrong-root trap is \(\lv3\); the coded centre-rooted pairing
(1800G) and interned-id counting are \(\lv4\); the hash-set merge (601D) is
\(\lv{5}\); the Summit (CF 1252F, every component the same shape) is step~6.}
\end{figure}

% ---------------- The running example ----------------
\subsection*{The Running Example: One Tree, Its Twin, and Two Traps}

Every deterministic section runs on the same rooted tree \(R\) on \(n=9\) nodes,
fixed once here and authoritative (if a later trace disagrees, this figure
wins).  Rooted at \(1\); parent-to-child edges, 1-based:
\[
  1\!\to\!2,\quad 1\!\to\!3,\quad 1\!\to\!4,\quad
  2\!\to\!5,\quad 2\!\to\!6,\quad 3\!\to\!7,\quad 3\!\to\!8,\quad 4\!\to\!9 .
\]
The root has \emph{three children, and two of them share a shape}: subtrees at
\(2\) and \(3\) are both ``a node above two leaves'' (isomorphic), while
subtree\((4)\) is ``a node above one leaf'' -- isomorphic to neither.  That is
the point: the pairing question at the root -- which children come
in mirror pairs, and may a middle child stand alone? -- drives
\(\secref{c22-centre}\) and the CF 1800G solve.  Depths are
\(\depth(1){=}0\); \(\depth(2){=}\depth(3){=}\depth(4){=}1\);
\(\depth(5){=}\dots{=}\depth(9){=}2\).  As an \emph{unrooted} tree, \(R\) has
diameter \(4\) (path \(5\)-\(2\)-\(1\)-\(3\)-\(7\)) and a \textbf{single centre}:
vertex \(1\).

\begin{figure}[H]
\centering
\begin{tikzpicture}[diagram, scale=0.95, transform shape,
   tnode/.style={circle, draw=primaryblue!70, fill=primaryblue!8, line width=0.7pt,
                 minimum size=6.6mm, font=\scriptsize\bfseries, inner sep=1pt},
   sub/.style={circle, draw=difforange, fill=difforange!16, line width=1.0pt,
             minimum size=6.6mm, font=\scriptsize\bfseries, inner sep=1pt}]
  \node[tnode] (n1) at (0,0) {1};
  \node[sub]   (n2) at (-3.2,-1.5) {2};
  \node[sub]   (n3) at (0,-1.5) {3};
  \node[tnode] (n4) at (3.2,-1.5) {4};
  \node[sub]   (n5) at (-4.0,-3.0) {5};
  \node[sub]   (n6) at (-2.4,-3.0) {6};
  \node[sub]   (n7) at (-0.8,-3.0) {7};
  \node[sub]   (n8) at (0.8,-3.0) {8};
  \node[tnode] (n9) at (3.2,-3.0) {9};
  \draw (n1)--(n2) (n1)--(n3) (n1)--(n4) (n2)--(n5) (n2)--(n6) (n3)--(n7) (n3)--(n8) (n4)--(n9);
  \node[annot, font=\scriptsize, anchor=north west, text width=8cm] at (-4.9,0.9)
    {orange = subtree(2) and subtree(3): the SAME shape;  node 4's subtree: a different one};
\end{tikzpicture}
\caption{The running example \(R\) (9 nodes, root 1).  The two orange subtrees
are isomorphic to each other (each is ``a node above two leaves''); subtree\((4)\)
(a node above one leaf) is isomorphic to neither.  Mirror pairing at the root and
the AHU ids of all nine subtrees are traced across 28.1 and 28.3.}
\label{fig:c22-tree}
\end{figure}

The \emph{unrooted twin} \(R'\) is \(R\) with every vertex renamed by
\(\sigma=(1\;2\;5)\) (figure
\(\figref{c22-twin}\) in \(\secref{c22-ahu}\)); vertex \(1\) of \(R'\) is
therefore vertex \(5\) of \(R\) -- a leaf, now wearing the root's hat.  \(R\)
and \(R'\) are isomorphic \emph{as unrooted trees} but not isomorphic \emph{as
rooted trees at the given roots} -- the one-sentence distinction that makes the
centre the canonical root, and the running thread of 28.3.

Two small traps return throughout.  The \textbf{look-alike trap}: the 7-vertex
trees \(T_1\) (edges \(1\!-\!2,2\!-\!3,3\!-\!4,4\!-\!5,5\!-\!6,3\!-\!7\)) and
\(T_2\) (edges \(1\!-\!2,2\!-\!3,3\!-\!4,4\!-\!5,5\!-\!6,2\!-\!7\)) -- the same
path of six with one pendant moved -- share the degree sequence
\((3,2,2,2,1,1,1)\), the diameter \(5\) and even the \emph{same centre edge}
\(\{3,4\}\), yet they are \textbf{not} isomorphic: each has exactly one
vertex of degree \(3\), and an isomorphism must map it to the other's while
preserving distances to the leaves -- \(T_1\)'s branch vertex sits at distances
\(\{1,2,3\}\) from the three leaves, \(T_2\)'s at \(\{1,1,4\}\).  Every cheap
invariant matches; only a canonical form separates them
(\(\secref{c22-centre}\)).  The \textbf{wrong-root trap} is the opposite
failure: \(R\) and \(R'\) above are the \emph{same} unrooted tree, yet their
rooted names at vertex \(1\) differ (\(4\) against \(7\)) -- rooting at an
arbitrary vertex answers a question nobody asked.  One move cures both: root at
the centre (Ch.~12's black box).  And when the centre is an \emph{edge} -- the
path \(P_6\) on vertices \(1\!-\!2\!-\!3\!-\!4\!-\!5\!-\!6\) has diameter \(5\)
and centres \(\{3,4\}\) -- ``root at the centre'' is replaced by the half-pairing
test of \(\secref{c22-centre}\).

\begin{figure}[H]
\centering
\begin{tikzpicture}[diagram, scale=0.9, transform shape,
   tnode/.style={circle, draw=primaryblue!70, fill=primaryblue!8, line width=0.7pt,
                 minimum size=6mm, font=\scriptsize\bfseries, inner sep=1pt},
   warn/.style={circle, draw=warningred, fill=warningred!12, line width=1pt,
                 minimum size=6mm, font=\scriptsize\bfseries, inner sep=1pt}]
  \node[tnode] (a1) at (0,0) {1};
  \node[tnode] (a2) at (1.15,0) {2};
  \node[warn]  (a3) at (2.3,0) {3};
  \node[tnode] (a4) at (3.45,0) {4};
  \node[tnode] (a5) at (4.6,0) {5};
  \node[tnode] (a6) at (5.75,0) {6};
  \node[tnode] (a7) at (2.3,-1.35) {7};
  \draw (a1)--(a2)--(a3)--(a4)--(a5)--(a6) (a3)--(a7);
  \node[annot, font=\scriptsize, anchor=south, text width=6.9cm, align=center] at (2.9,0.42)
    {$T_1$: pendant on vertex 3; the branch vertex sees the leaves at 1, 2, 3};
  \node[tnode] (b1) at (0,-3.1) {1};
  \node[warn]  (b2) at (1.15,-3.1) {2};
  \node[tnode] (b3) at (2.3,-3.1) {3};
  \node[tnode] (b4) at (3.45,-3.1) {4};
  \node[tnode] (b5) at (4.6,-3.1) {5};
  \node[tnode] (b6) at (5.75,-3.1) {6};
  \node[tnode] (b7) at (1.15,-4.45) {7};
  \draw (b1)--(b2)--(b3)--(b4)--(b5)--(b6) (b2)--(b7);
  \node[annot, font=\scriptsize, anchor=north, text width=6.9cm, align=center] at (2.9,-3.52)
    {$T_2$: pendant on vertex 2; the branch vertex sees the leaves at 1, 1, 4};
  \node[annot, font=\scriptsize, text=warningred, anchor=west] at (6.1,-0.25) {centre edge 3-4};
  \node[annot, font=\scriptsize, text=warningred, anchor=west] at (6.1,-3.35) {centre edge 3-4};
\end{tikzpicture}
\caption{The look-alike pair \((T_1,T_2\)): same invariants, different trees.
Every cheap invariant says yes: same size, same degree sequence
\((3,2,2,2,1,1,1)\), same diameter \(5\), same centre edge \(\{3,4\}\) (red
labels mark the unique degree-\(3\) vertex).  The decisive read is the distance
profile from that vertex to the leaves, \(\{1,2,3\}\) against \(\{1,1,4\}\) --
isomorphisms preserve distances, so no bijection exists.  \(\secref{c22-centre}\)
computes the same verdict mechanically from centre-rooted canonical names, and
that machinery also catches the opposite error (\(R\) versus \(R'\): isomorphic,
yet wrong-rooted names disagree).}
\label{fig:c22-trap}
\end{figure}

% ---------------- Which job, cost chart ----------------
\subsection*{Which Job? (the question-shape chart)}

\begin{center}
\small
\renewcommand{\arraystretch}{1.25}
\setlength{\tabcolsep}{3pt}
\begin{tabular}{@{}p{0.31\linewidth} p{0.19\linewidth} p{0.27\linewidth} >{\raggedright\arraybackslash}p{0.15\linewidth}@{}}
\toprule
\textbf{Question} & \textbf{Technique} & \textbf{Why plain per-node DP misses it} & \textbf{Where} \\
\midrule
same rooted shape as that subtree? & AHU interned ids & the DP folds \emph{values}; here the folded object is the shape itself & \secref{c22-ahu} \\
children mirror-pairable at every node? & canonical hash + pairing & "equal shape" is the predicate; no arithmetic statistic captures it & \secref{c22-centre} \\
\(O(1)\) shape compares after \(O(n)\) & subtree hashing & re-running AHU per compare is \(\bigO{n}\); a hash trades certainty for speed & \secref{c22-hash} \\
how many \emph{different} downward strings? & hash sets + small-to-large & strings are compared as hashes; sets merged by Ch.17's union & \secref{c22-hash} \\
count distinct subtree shapes & interned id multiplicity & the count is over \emph{isomorphism classes}, not vertices & \secref{c22-distinct} \\
cut one node -- all pieces identical? & all-hashes + reroot & the parent-side piece needs the reroot view, not the child view & \(\probref{c22-summit}\) \\
\bottomrule
\end{tabular}
\end{center}

\begin{insightbox}[One promise decides everything: sort before you compare]
Every tool in this chapter is the same idea at a different intensity.  Two
unordered collections are equal iff their \emph{sorted} forms are equal -- so
make the representation order-free \textbf{before} comparing: AHU sorts child
labels and interns the tuple; a hash combines child values
\emph{commutatively} (XOR/add, so insertion order cannot matter); the unrooted
case sorts by choosing a canonical root -- the centre.  The whole chapter is one
sentence: \emph{remove every accidental degree of freedom (label names, child
order, root choice) and then equality is just equality.}  What is left after
you remove those three is the \textbf{shape} -- the object this chapter computes.
\end{insightbox}

% ---------------- Taxonomy chart ----------------
\begin{figure}[H]
\centering
\begin{tikzpicture}[diagram, scale=0.95, transform shape,
    mnode/.style={circle, draw=#1, line width=0.9pt, fill=#1!8, minimum size=9mm,
                  font=\scriptsize\bfseries, text=#1!60!black}]
  \node[mnode=primaryblue] (a) at (0,0)   {intern};
  \node[annot, font=\scriptsize, anchor=west, text width=9.8cm, align=left] at (1.7,0) {rooted AHU: label leaves first, sort each node's child labels, give each fresh sorted tuple a new id; root ids equal = isomorphic. O(n log n) (28.1)};
  \node[mnode=successgreen] (b) at (0,-1.1) {center};
  \node[annot, font=\scriptsize, anchor=west, text width=9.8cm, align=left] at (1.7,-1.1) {unrooted: root at the one center (Ch.12 black box) and compare canonical forms; two centers = pair the halves at the center edge -- rooting anywhere else is the wrong-root trap (28.3)};
  \node[mnode=difforange]   (c) at (0,-2.2) {hash};
  \node[annot, font=\scriptsize, anchor=west, text width=9.8cm, align=left] at (1.7,-2.2) {randomize for O(1) compares: h(u) = commutative fold (sum/XOR) of random per-child, per-depth values; 64-bit or double modulus; keep a collision ledger and re-verify YES by an explicit bijection DFS (28.2)};
  \node[mnode=softpurple]   (d) at (0,-3.3) {count};
  \node[annot, font=\scriptsize, anchor=west, text width=9.8cm, align=left] at (1.7,-3.3) {interned ids are shape names: a map id -> multiplicity counts isomorphism classes of subtrees; sets of downward-string hashes merged small-to-large count distinct downward strings (28.4, 601D)};
  \node[mnode=warningred]   (e) at (0,-4.4) {reject};
  \node[annot, font=\scriptsize, anchor=west, text width=9.9cm, align=left] at (1.7,-4.4) {labels matter -> tree DP / suffix tools; order matters -> ordered-tree encodings; online edits -> HLD (Ch.22) / LCT (Ch.25); ancestor set batches -> Mo (Ch.27); distance distributions -> centroid (Ch.24) (28.5)};
\end{tikzpicture}
\caption{The chapter's five-state taxonomy.  \emph{intern} is the deterministic
AHU fold; \emph{center} lifts it to unrooted trees via the canonical root;
\emph{hash} trades determinism for \(O(1)\) compares; \emph{count} uses ids as
shape names; \emph{reject} hands the job to the right neighbour.  Each solved
problem sits on one row: CF 1800G on \emph{center}+pairing, CF 601D on
\emph{hash}+count, the Summit CF 1252F on \emph{hash}+center at the edge.}
\label{fig:c22-taxonomy}
\end{figure}

% ---------------- Who owns what ----------------
\subsection*{Who Owns What (so nothing is learned twice)}

\begin{center}
\small
\renewcommand{\arraystretch}{1.25}
\setlength{\tabcolsep}{3pt}
\begin{tabular}{@{}p{0.28\linewidth} p{0.24\linewidth} p{0.44\linewidth}@{}}
\toprule
\textbf{Idea} & \textbf{Owner} & \textbf{This chapter's share} \\
\midrule
diameter, centre, eccentricity & Ch.~12 (\chref{c09}) & the centre is \emph{used} as the canonical root; its computation is not re-derived \\
the bottom-up fold & Ch.~13 (\chref{c10}) & the same recursion, folding \emph{shapes} (tuples/ids) instead of numbers \\
small-to-large set merging & Ch.~17 (\chref{c28}) & the 601D hash-set unions reuse its complexity argument verbatim \\
rerooting / parent-side values & Ch.~15 (\chref{c20}) & the Summit's parent-side component hash is the standard reroot template \\
string hashing (poly mod) & the string chapters & 28.2 needs a different animal: an \emph{order-free} (commutative) combine \\
``which subtree equals which'' & Ch.~27 boundary (\chref{c21}) & the baton pass -- 28.1--28.4 answer it in \(O(n\log n)\) total \\
\bottomrule
\end{tabular}
\end{center}

% ---------------- Recognition guide ----------------
\begin{recognitionbox}[Pattern card -- ``same shape, ignore the labels'']
\reckey{statements where two trees (or two subtrees of one tree) must be called
equal or counted \emph{up to relabelling}: ``identical'', ``the same structure'',
``mirror image'', ``can the children be ordered so that \(\dots\)'', ``pairwise
same trees after a cut''.  The tell is that \textbf{vertex labels and numbers
carry no information} -- only adjacency matters -- and usually child order is
free too.  Variants that keep the tell: counting how many \emph{distinct shapes}
occur; finding the vertex whose incident pieces agree.}
\recuse{
  \item rooted shape equality / mirror pairing -- canonical ids by AHU (\(\secref{c22-ahu}\), \(\secref{c22-centre}\), solved CF 1800G);
  \item many repeated compares inside one tree -- subtree hashes, \(O(1)\) each (\(\secref{c22-hash}\));
  \item distinct downward strings per node -- hash sets + small-to-large (\(\secref{c22-hash}\), solved CF 601D);
  \item isomorphism-class counts -- interned id multiplicities (\(\secref{c22-distinct}\));
  \item every component of a cut identical -- all-hashes + reroot (\(\probref{c22-summit}\)).
}
\recnotuse{
  \item vertex \emph{values/colours} must match too -- plain tree DP (\chref{c10}) or DSU-on-tree (\chref{c28}); shape-only tools must then compare \((\text{shape},\text{value})\) pairs;
  \item strings read along paths with labels mattering -- suffix automaton / hashing on chains, not shape hashes;
  \item the question is about \emph{distances} -- Ch.~12 metrics or centroid (\chref{c18});
  \item child order is fixed (ordered trees) -- skip the sort, an in-order encoding suffices;
  \item trees change online -- link-cut (\chref{c23}) territory.
}
One-line test: \emph{if renaming every vertex (and reordering every child list)
cannot change the answer, you are comparing shapes -- canonical form or
commutative hash; if a name or an order can change the answer, it is not this
chapter.}
\end{recognitionbox}

% ---------------- Constraint -> technique ----------------
\subsection*{Reading the Constraints (which technique the limits want)}

\begin{center}
\small
\renewcommand{\arraystretch}{1.25}
\setlength{\tabcolsep}{3pt}
\begin{tabular}{@{}p{0.34\linewidth} p{0.40\linewidth} p{0.20\linewidth}@{}}
\toprule
\textbf{Constraint signature} & \textbf{What it rules in/out} & \textbf{Technique} \\
\midrule
compare two trees once, \(n\le2\cdot10^5\) & one canonical form each, compare roots & \secref{c22-ahu} \\
\(t\le10^4\) testcases, \(\sum n\le2\cdot10^5\), YES/NO per tree & ids must be re-fresh per case or use hashing & \secref{c22-centre} (1800G) \\
thousands of subtree-vs-subtree compares inside one tree & AHU per compare is too slow; precompute hashes & \secref{c22-hash} \\
\(n\le3\cdot10^5\), distinct strings per node & sets + small-to-large, 64-bit hashes & \secref{c22-hash} (601D) \\
\(n\le4000\), ``does ANY cut split into identical pieces'' & \(O(n^2)\) compares legal; prune candidates first & \(\probref{c22-summit}\) (1252F) \\
\bottomrule
\end{tabular}
\end{center}

% ---------------- Notation ----------------
\subsection*{Notation Used Only Here}

The fixed running tree is \(R\), rooted at \(1\); \(T_1,T_2\) are the look-alike pair
and \(P_6\) the two-centre path, all defined above.  \(\children(u)\) is the
multiset of children of \(u\) in the rooted view; \(\subtree(u)\) is the rooted
subtree at \(u\).  A \emph{canonical label} \(\mathrm{cl}(u)\) is the interned id of the
sorted tuple of \(\mathrm{cl}\) values of \(u\)'s children (28.1);
\(\Phi(T)\) is the canonical form of a whole (un)rooted tree \(T\).  \(h(u)\) is
the random \emph{subtree hash} of \(\subtree(u)\) (28.2); \(\mathcal{S}(v)\) is
the set of hashes of strings read downward from \(v\) (601D).  \(\mathrm{centre}(T)\) is
the centre set of an unrooted tree (one vertex or one edge, Ch.12).  ``\(\cong\)''
reads ``is isomorphic to''; a superscript \(^{\ast}\) marks a verdict that is
only \(O(1)\)-certain \emph{up to the collision ledger} of \(\secref{c22-hash}\).
\(\depth,\size\) are Ch.~5/Ch.~12's, used unchanged.  Every section runs on
\((R,R',T_1,T_2,P_6)\) above; if a table here and a trace there disagree, this
opener is authoritative.